\medskip

\textit{Pour chacune des affirmations suivantes, indiquer si elle est vraie ou fausse. Chaque réponse doit être justifiée. Une réponse non justifiée ne rapporte aucun point.}

\smallskip

\textit{Les cinq questions de cet exercice sont indépendantes.}

\begin{enumerate}
\item On considère une suite $\left(t_n\right)$ vérifiant la relation de récurrence : \[ \text{pour tout entier naturel } n,~t_{n+1}= - 0,8t_n + 18.\]

	\textbf{Affirmation 1} : La suite $\left(w_n\right)$ définie pour tout entier naturel $n$ par $w_n = t_n - 10$ est géométrique.
	
\smallskip

\item On considère une suite $\left(S_n\right)$ qui vérifie pour tout entier naturel $n$ non nul : \[ 3n - 4 \leqslant S_n \leqslant 3n + 4.\]
	
	La suite $\left(u_n\right)$ est définie, pour tout entier naturel $n$ non nul, par : $u_n=\frac{S_n}{n}$.
	
	\textbf{Affirmation 2} : La suite $\left(u_n\right)$ converge.
	
	\smallskip
\item On considère la suite $\left(v_n\right)$ définie par :
	\[ v_1=2 \text{ et pour tout entier naturel } n\geqslant1,~v_{n+1}=2-\frac{1}{v_n}. \]
	
	\textbf{Affirmation 3} : Pour tout entier naturel $n \geqslant 1$, $v_{n}=\dfrac{n+1}{n}$.
	
	\smallskip
\item On considère la suite $\left(u_n\right)$ définie pour tout entier naturel $n$ par $u_n = \e^n-n$.
	
	\textbf{Affirmation 4} : La suite $\left(u_n\right)$ converge.
	
	\smallskip
\item On considère la suite $\left(u_n\right)$ définie à l'aide du script écrit ci-dessous en langage \textsf{Python}, qui renvoie la valeur de $u_n$.
	
\texttt{\begin{center}
\begin{tabular}{|l|}\hline
\textbf{def} u(n) :\\
\qquad valeur = 2\\
\qquad \textbf{for} k \textbf{in range}(n) :\\
\qquad \qquad valeur = 0.5 * (valeur + 2/valeur)\\
\qquad \textbf{return} valeur\\ \hline
\end{tabular}
\end{center}
}

On admet que $\left(u_n\right)$ est décroissante et vérifie pour tout entier naturel $n$ :

\[\sqrt{2} \leqslant u_n \leqslant 2.\]

\textbf{Affirmation 5} : La suite $\left(u_n\right)$ converge vers $\sqrt{2}$.
\end{enumerate}

\bigskip
